\chapter{Discussion}
\label{ch:discussion}

\section{What worked, and what did not}
The clearest result of this project is negative, and it repeated itself. Several
choices that looked as though they should help turned out not to, and in each
case this became visible only through a controlled test.

\begin{itemize}
  \item \tfunet{}'s poor first score looked like a weakness of its architecture.
    It came instead from two training settings, class weighting and per-image
    normalisation; correcting them took it from 0.39 to 0.93 without changing
    the network (\cref{ch:tfunet}).
  \item The components \hybrid{} was designed around---strip convolutions,
    channel attention and residual blocks---contribute nothing measurable on
    real data, not even for the faint RFI they were meant to target
    (\cref{ch:drivers}).
  \item A CRF refinement that works well on top of a simple threshold added
    nothing on top of a network that already uses neighbourhood information,
    and with its published settings it was harmful (\cref{sec:crf}).
\end{itemize}

The model still matches or slightly exceeds the best published results on
LOFAR, but the reason is not the one it was designed around. Without the ablation, I would have
credited the wrong parts of the model.

\section{Synthetic benchmarks can mislead}
The same network, trained the same way, scored 0.93 on my synthetic data and
0.49 on LOFAR. The synthetic labels are noise-free and every labelled pixel is
above the noise by construction, while on real data a fifth of the RFI is
fainter than a typical clean pixel and 44\% of the pixels flagged in the
training labels are not RFI according to a human expert. The synthetic data was useful for finding
and fixing \tfunet{}'s training problems quickly, but it was a poor guide to
how methods compare on real observations: on synthetic data every variant of
my model scored above 0.95, so it could not show which parts mattered.

\section{Single runs can mislead}
Three times in this project, an effect measured with one training run shrank
or disappeared when the run was repeated with different random seeds:
\begin{itemize}
  \item lowering \tfunet{}'s learning rate first appeared to add 0.118 to
    \Fone{}; over three seeds the mean gain was 0.058 and not significant
    (\cref{sec:tfunet_lofar});
  \item removing GroupNorm first appeared to cost 0.034; over three seeds the
    difference was 0.009 and not significant (\cref{sec:ablation});
  \item strip convolutions first appeared to help faint RFI by about 0.02; the
    effect came from comparing models at different thresholds and vanished
    when they were compared fairly (\cref{sec:faint}).
\end{itemize}
On LOFAR, results vary between seeds by up to about 0.05 in \Fone{}, which is as
large as many of the effects one might want to measure. Repeating runs and
comparing models under identical conditions was the only reliable way to tell
real effects from noise.

\section{Limitations}
\begin{enumerate}
  \item \textbf{Published baselines are quoted.} The results for
    Mesarcik et al.'s U-Net and RFI-Net, RFDL and Swin-UNETR are taken from
    their papers\footcite{mesarcik2022,vanzyl2024rfdl,ouyang2024swin}. Because my
    measurement of \aoflagger{} matches the value these papers report to four
    decimal places, the test data, labels and metric are the same, so the
    comparison is like for like; the values themselves are as they reported
    them.
  \item \textbf{Small test set.} There are only 109 human-labelled test images.
    The resulting uncertainty, about $\pm0.04$ in \Fone{} at 95\% confidence, is
    larger than the variation between seeds.
  \item \textbf{Different training budgets.} Mesarcik et al.\ trained on
    32$\times$32 patches for 100 epochs; I trained on full images for a fixed
    number of steps. Each method was trained in a sensible way for it, but not
    under one common budget.
  \item \textbf{Three seeds.} This is the minimum for any statistical
    statement, and several conclusions rest on it.
  \item \textbf{One real dataset.} All real-data results come from LOFAR. The
    method has not been tested on data from another telescope, including the
    GMRT data that originally motivated this project.
  \item \textbf{The gain over \tfunet{} is not yet explained}
    (\cref{sec:unexplained}).
  \item \textbf{\tfunet{} was not run on HERA}, so it is compared with my model
    only on the synthetic data and on LOFAR.
\end{enumerate}

\section{Future work}
The results suggest that further additions to the model should target a
specific, measured weakness, and each should be tested by removing it again.
\begin{itemize}
  \item \textbf{Find the source of the current gain.} Train the plain U-Net
    without the Dice loss, then test the other remaining differences from
    \tfunet{} one at a time. This is cheap and would explain the main result.
  \item \textbf{Errors at the edges of RFI.} Missed RFI is the larger source
    of error: removing every missed pixel would raise \Fone{} from about 0.66 to
    0.82, and removing every false alarm would raise it to 0.77. In one model
    examined in detail, 56\% of the missed pixels lie within two pixels of a
    correct detection and 77\% of the false alarms within two pixels of real
    RFI, so most errors concern the exact outline of an RFI feature rather than
    whole features.
  \item \textbf{Learn from noisy labels.} Every training label comes from
    \aoflagger{}, and 44\% of the pixels it flags are not RFI according to the
    human expert. Two simple remedies, softening or ignoring the one-pixel edge
    of every \aoflagger{} mask during training, were tried after the main
    experiments; both lowered \Fone{} (by 0.008 and 0.011, significant over
    three seeds).
  \item \textbf{Use a few expert labels.} In a pilot, fine-tuning on 20
    expert-labelled images raised \Fone{} on 70 other expert-labelled images by
    0.018, in all 9 runs, while the same fine-tuning with \aoflagger{}'s labels
    gave no gain. A cross-validated study that also chooses which image patches
    to label is the next step.
  \item \textbf{Faint RFI.} Only about a quarter of the faint and invisible RFI
    is found. Since these pixels carry no brightness signal of their own, any
    improvement has to come from wider context.
  \item \textbf{Apply the method to GMRT pulsar data}, the original motivation
    for this project.
\end{itemize}
